\documentclass[9pt,a4paper]{extarticle}
\usepackage{fontspec}
\usepackage[greek,english]{babel}
\babelfont{rm}{Noto Serif}
\babelfont{sf}{Noto Sans}
\setsansfont{Noto Sans}
\setmainfont{Noto Serif}
\usepackage[a4paper,margin=1.55cm,top=1.35cm,bottom=1.35cm]{geometry}
\usepackage{microtype}
\usepackage{amsmath,amssymb,mathtools}
\usepackage{booktabs,tabularx,array,longtable}
\usepackage{enumitem}
\usepackage{xcolor}
\usepackage[most]{tcolorbox}
\usepackage{fancyhdr}
\usepackage{hyperref}
\usepackage{ragged2e}
\usepackage{titlesec}
\usepackage{setspace}
\usepackage{pifont}
\hypersetup{colorlinks=true,linkcolor=black,urlcolor=black}
\definecolor{navy}{HTML}{18324A}
\definecolor{slate}{HTML}{4B5D67}
\definecolor{light}{HTML}{F3F6F8}
\definecolor{warm}{HTML}{F8F3EC}
\definecolor{greenbox}{HTML}{EEF6F0}
\definecolor{redbox}{HTML}{FAEEEE}
\definecolor{bluebox}{HTML}{EEF3F8}
\setlength{\parindent}{0pt}
\setlength{\parskip}{4pt}
\emergencystretch=4em
\sloppy
\setlist[itemize]{leftmargin=1.35em,itemsep=1.5pt,topsep=2pt}
\setlist[enumerate]{leftmargin=1.45em,itemsep=1.5pt,topsep=2pt}
\titleformat{\section}{\sffamily\bfseries\Large\color{navy}}{}{0pt}{}
\titleformat{\subsection}{\sffamily\bfseries\normalsize\color{slate}}{}{0pt}{}
\pagestyle{fancy}
\fancyhf{}
\fancyhead[L]{\sffamily\small Εγχειρίδιο μελέτης & έρευνας}
\fancyhead[R]{\sffamily\small v8.71a / M1--M5}
\fancyfoot[C]{\sffamily\small \thepage}
\newtcolorbox{keybox}[1][]{colback=bluebox,colframe=navy!65!black,boxrule=.5pt,arc=2mm,left=2mm,right=2mm,top=1.5mm,bottom=1.5mm,#1}
\newtcolorbox{warnbox}[1][]{colback=warm,colframe=slate!70!black,boxrule=.5pt,arc=2mm,left=2mm,right=2mm,top=1.5mm,bottom=1.5mm,#1}
\newtcolorbox{redline}[1][]{colback=redbox,colframe=red!45!black,boxrule=.5pt,arc=2mm,left=2mm,right=2mm,top=1.5mm,bottom=1.5mm,#1}
\newtcolorbox{greenline}[1][]{colback=greenbox,colframe=green!45!black,boxrule=.5pt,arc=2mm,left=2mm,right=2mm,top=1.5mm,bottom=1.5mm,#1}
\newcommand{\say}[1]{\begin{keybox}\textbf{Προφορική διατύπωση:} #1\end{keybox}}
\newcommand{\avoid}[1]{\begin{redline}\textbf{Μην το πεις έτσι:} #1\end{redline}}
\newcommand{\ask}[1]{\begin{warnbox}\textbf{Ερώτηση προς τον επιβλέποντα:} #1\end{warnbox}}
\newcommand{\formula}[1]{\begin{center}\fbox{\(#1\)}\end{center}}
\begin{document}
\selectlanguage{greek}

% PAGE 1
\begin{center}
{\sffamily\bfseries\LARGE Μικρογραφία της ερευνητικής πρότασης}\\[3mm]
{\sffamily\Large Τα απολύτως απαραίτητα: έννοιες, μεθοδολογία, απαντήσεις, όρια ισχυρισμών}\\[2mm]
{\large «Η μουσική θέση ως υπολογιστική μονάδα της βυζαντινής μουσικής»}\\[4mm]
{\small Βάση: αναθεωρημένη πρόταση v8.71a, κανόνας επιλογής \(R^\star\), EISE, M3--M5}
\end{center}

\begin{keybox}
\textbf{Η διατριβή σε μία πρόταση.} Ελέγχω αν η μουσική θέση μπορεί να τυποποιηθεί ως μεταβλητού μήκους, μουσικολογικά ερμηνεύσιμη και υπολογιστικά ελέγξιμη ενδιάμεση μονάδα, ώστε να αναγνωρίζονται οικογένειες, να μελετάται η ιστορικά τεκμηριωμένη μεταβολή τους, να συνδέονται με ακρόαση και εκτέλεση και τελικά να χαρακτηρίζονται πλήρη έργα ως προς δηλωμένο σώμα αναφοράς.
\end{keybox}

\subsection*{Ο χάρτης που πρέπει να έχεις στο μυαλό σου}
\begin{center}
\begin{tabularx}{.95\linewidth}{>{\bfseries\sffamily}p{1.25cm}X}
M1 & Ποια αναπαράσταση διατηρεί καλύτερα τα όρια και την οικογενειακή ταυτότητα; Επιλογή \(R^\star\).\\
M2 & Προσθέτουν οι ιστορικά τεκμηριωμένες εξηγητικές σχέσεις πληροφορία πέρα από απλή ενδοοικογενειακή συγγένεια; EISE έναντι B2.\\
M3a & Είναι οι οικογένειες και τα όρια μουσικολογικά έγκυρα σε ανεξάρτητη αξιολόγηση;\\
M3c & Προσθέτει η οικογενειακή προσδοκία πληροφορία στην ανθρώπινη κρίση πέρα από την event-level προσδοκία;\\
M3b & Οργανώνουν οι ακροατές τις πραγματώσεις κατηγορικά σύμφωνα με τις οικογένειες πέρα από επιφανειακή ομοιότητα;\\
M4 & Συμφωνεί η ιστορικά οργανωμένη παραλλαγή με ανεξάρτητη σύγχρονη εκτελεστική μεταβλητότητα;\\
M5 & Πώς διαφοροποιείται ένα πλήρες έργο ως προς λεξιλόγιο, διαδοχή, λειτουργική τοποθέτηση και εσωτερική πραγμάτωση;\\
\end{tabularx}
\end{center}

\say{«Δεν ξεκινώ από την παραδοχή ότι η θέση είναι εκ των προτέρων η καλύτερη computational unit. Την κάνω ελέγξιμη υπόθεση, συγκρίνω αναπαραστάσεις στη M1, και μόνο η οριστικοποιημένη \(R^\star\) περνά στην κύρια M2.»}

\begin{warnbox}
\textbf{Στόχος προετοιμασίας.} Δεν χρειάζεται να αποδείξεις ότι όλα έχουν ήδη υλοποιηθεί. Χρειάζεται να δείξεις ότι γνωρίζεις: ποιο είναι το επιστημονικό ερώτημα, γιατί κάθε μέθοδος υπάρχει, ποιο αποτέλεσμα θα άλλαζε την ερμηνεία σου, ποια σημεία χρειάζονται τεχνική καθοδήγηση και ποιο είναι το ελάχιστο βιώσιμο PhD.
\end{warnbox}

\newpage
% PAGE 2
\section*{1. Πρόβλημα, κενό και συμβολή}

\subsection*{Το πρόβλημα}
Η Νέα Μέθοδος δίνει πλούσια σημειογραφική επιφάνεια, αλλά το επιστημονικό αντικείμενο της διατριβής βρίσκεται σε ενδιάμεση κλίμακα: ούτε μεμονωμένο νεύμα ούτε ολόκληρη φράση. Η μουσικολογική παράδοση χρησιμοποιεί την έννοια της \textbf{θέσης} για επαναχρησιμοποιούμενες, λειτουργικά ενταγμένες μελωδικές δομές. Η δυσκολία είναι ότι οι θέσεις έχουν μεταβλητό μήκος, μεταβάλλονται αναλυτικά, έχουν αβέβαια όρια και εμφανίζονται σε διαφορετικά τροπικά, κειμενικά και ιστορικά συμφραζόμενα.

\subsection*{Το ερευνητικό κενό}
Υπάρχουν υπολογιστικές προσεγγίσεις για συμβολική αναπαράσταση, pattern discovery, segmentation, similarity, chant modelling και prediction. Υπάρχει επίσης πλούσια μουσικολογική θεωρία της βυζαντινής μουσικής για θέσεις, εξηγήσεις, τροπικότητα και μορφολογία. Το κενό είναι η \textbf{ελεγχόμενη σύνδεση των δύο}: να μετατραπεί η θέση σε computational object χωρίς να χαθεί η ιστορική/μουσικολογική της σημασία και χωρίς να θεωρηθεί ground truth επειδή απλώς υπάρχει στη θεωρία.

\subsection*{Η κύρια συμβολή}
\begin{itemize}
\item Τυποποίηση της θέσης ως μεταβλητού μήκους μονάδας με ρητή provenance και αβεβαιότητα.
\item Πειραματική επιλογή αναπαράστασης αντί αυθαίρετης tokenization.
\item B2--EISE σχεδιασμός που απομονώνει την \textbf{προέλευση} των θετικών σχέσεων.
\item Σύνδεση συμβολικής ιστορικής παραλλαγής με ανεξάρτητη εκτελεστική μεταβλητότητα.
\item Ανθρώπινη επικύρωση κατηγορικής οργάνωσης και προσδοκίας.
\item Εφαρμογή σε πλήρη έργα με πολυδιάστατο προφίλ και επιστροφή από κάθε αριθμητικό εύρημα στο συγκεκριμένο μουσικό χωρίο.
\end{itemize}

\subsection*{HarmonyTok: τι ακριβώς κρατάς}
Το HarmonyTok είναι \textbf{methodological precedent}: η tokenization δεν θεωρείται ουδέτερη προεπεξεργασία αλλά πειραματική μεταβλητή που συγκρίνεται σε κοινά tasks. Δεν αποτελεί μοντέλο της βυζαντινής θέσης και δεν θεμελιώνει τον EISE.

\say{«Το ερευνητικό κενό δεν είναι ότι δεν έχει γίνει machine learning στη βυζαντινή μουσική. Είναι ότι η ιστορικά νοηματοδοτημένη θέση δεν έχει ακόμη ελεγχθεί ως computational unit με αυστηρή διάκριση representation, segmentation, family identity, historical relation και external validation.»}

\ask{«Θεωρείτε ότι η κύρια συνεισφορά πρέπει να παρουσιαστεί πρωτίστως ως MIR representation problem ή ως computational musicology problem με MIR μεθοδολογία; Θα άλλαζε αυτό τον τρόπο που θα δομούσατε τη M1;»}

\newpage
% PAGE 3
\section*{2. Μουσικολογική βάση: πέντε έννοιες, τρεις σχέσεις}
\subsection*{Η βασική οντολογία}
\begin{tabularx}{\linewidth}{>{\bfseries}p{3cm}X}
Θέση (τύπος) & Μουσικός σχηματισμός με διαγνωστικές σχέσεις, μεταβλητή έκταση και τροπική/φραστική λειτουργία.\tabularnewline
Οικογένεια & Επιχειρησιακή αναλυτική κατηγορία συγγενών πραγματώσεων· δεν είναι αυτόματα ιστορική ονομασία.\tabularnewline
Εμφάνιση & Συγκεκριμένο τμήμα ενός έργου με όρια, προέλευση και ενδεχόμενη αβεβαιότητα.\tabularnewline
Πραγμάτωση & Η συγκεκριμένη σημειογραφική/μελωδική μορφή που έχει μια θέση σε αυτή την εμφάνιση.\tabularnewline
Ιστορική εξήγηση & Σχέση συγκεκριμένων πηγών/αποδόσεων με αποδείξιμη ιστορική προέλευση· όχι απλή ομοιότητα.\tabularnewline
\end{tabularx}

\subsection*{Τρία διαφορετικά είδη σχέσεων}
\[E_{sim}\neq E_{fam}\neq E_{ex}.\]
Η μελωδική ομοιότητα δεν αρκεί για οικογενειακή ένταξη, και η ίδια οικογένεια δεν σημαίνει ιστορικά τεκμηριωμένη εξηγητική σχέση. Οι ακμές εξήγησης απαιτούν πηγή, κοινό χωρίο και έλεγχο τεκμηρίων.

\subsection*{Ιστορική θεμελίωση: τι αποδίδουμε με ασφάλεια}
Ο Μανουήλ Χρυσάφης και ο Γαβριήλ ιερομόναχος αποτελούν ιστορικά σημεία αναφοράς για τη θέση ως οργανωμένο μελικό σχηματισμό, όχι για τη συγκεκριμένη σημερινή υπολογιστική οικογένεια. Στον Χρύσανθο, ειδικά στο \emph{Θεωρητικόν Μέγα} §408, η εξάρτηση της ενέργειας παλαιότερων σχηματισμών από τη βαθμίδα εφαρμογής ενισχύει τη σημασία του συμφραζομένου· δεν αποτελεί προδιαγραφή του EISE ή του PPM. Η κριτική εξέταση των πηγών είναι αναγκαία όταν ισχυριζόμαστε εξηγητική σχέση.

\subsection*{Γιατί μεταβλητό μήκος;}
Η θέση δεν συμπίπτει υποχρεωτικά με ένα νεύμα, ένα ατομικό γεγονός, μία συλλαβή ή ολόκληρη φράση. Μπορεί να υπάρχουν συνδετικές και ανεπίλυτες περιοχές, καθώς και τεκμηριωμένες επικαλύψεις. Η βασική έξοδος δεν επιβάλλει ετικέτα οικογένειας σε κάθε γεγονός.

\say{«Ο EISE μαθαίνει συγγένεια πραγματώσεων από ελεγχόμενα δεδομένα· δεν βαφτίζει μια ομοιότητα ιστορική εξήγηση. Η οικογένεια είναι ελέγξιμη μουσικολογική υπόθεση και όχι δεδομένο αξίωμα.»}
\avoid{«Όποια δύο αποσπάσματα είναι συγγενή αποτελούν εξηγητικό ζεύγος» ή «ο Χρύσανθος θεμελιώνει το συγκεκριμένο PPM».}
\newpage
% PAGE 4
\section*{3. ΠΕ0: πηγές, κάρτες οικογενειών, επισημείωση}
\subsection*{Τα δεδομένα έχουν διαφορετικούς ρόλους}
\begin{itemize}
\item \textbf{Κατάλογος}: η παπαδική συλλογή της Μονής Αγίου Αντωνίου καλύπτει υλικό προέλευσης από 758 συνθέσεις· είναι καταλογογραφική αφετηρία, όχι αυτόματη βάση πλήρως επισημειωμένων θέσεων. Οι 10.000+ φόρμουλες ανήκουν σε διαφορετική, ειρμολογική/στιχηραρική συλλογή.
\item \textbf{Μητρώο εμφανίσεων}: έργο, πηγή, ήχος, κειμενικό χωρίο, όρια, λειτουργία, ανεπίλυτο υλικό, επιτρεπτές εναλλακτικές, προέλευση.
\item \textbf{Ιστορικός γράφος}: χωριστή σχέση εξήγησης, κοινό χωρίο, ανεξαρτησία και εξάρτηση πηγών. Δεν παράγουμε εξηγητικά positives συνθετικά.
\item \textbf{Ανεξάρτητη αξιολόγηση}: έργα και πηγές εκτός εκπαίδευσης/επιλογής, ώστε να μην υπάρχει διαρροή.
\end{itemize}
\subsection*{Κάρτα οικογένειας: τι ορίζουμε πριν μάθει το μοντέλο}
Αναλυτική κλίμακα, αναγκαίες μελωδικές σχέσεις, επιτρεπτές μεταβολές, τροπικοί περιορισμοί όπου τεκμηριώνονται, θετικά και δύσκολα αρνητικά παραδείγματα, κανόνες μη ένταξης και δυνατότητα \texttt{unresolved}. Οι πρώτες υποψήφιες περιπτώσεις (π.χ. παρακλητική, κύλισμα, πίεσμα) δεν θεωρούνται αυτόματα επικυρωμένες οικογένειες.
\subsection*{Μην συγχέεις τις τρεις φάσεις επισημείωσης}
\begin{tabularx}{\linewidth}{>{\bfseries}p{3cm}X}
ΠΕ0 / ανάπτυξη & Εσύ και δεύτερος ανεξάρτητος ειδικός σε προκαθορισμένο μέρος· δημιουργία και έλεγχος εγχειριδίου και δεδομένων ανάπτυξης.\tabularnewline
Πιλότος οδηγιών & 2--4 ενεργοί ψάλτες, χωριστοί από τη M3a· ελέγχουν εάν η εργασία είναι κατανοητή. Οι ετικέτες τους \emph{δεν} γίνονται επιβεβαιωτικά δεδομένα.\tabularnewline
M3a / τυφλή μελέτη & 3--5 άλλοι ανεξάρτητοι ειδικοί, σε χωριστό σύνολο, χωρίς προβλέψεις μοντέλων και χωρίς τις κρίσεις των άλλων.\tabularnewline
\end{tabularx}
\say{«Το ΠΕ0 δεν παράγει δήθεν έτοιμο gold standard: συγκροτεί μουσικολογικά ελέγξιμα δεδομένα με προέλευση, αβεβαιότητα και χωριστές συνθήκες ανάπτυξης και αξιολόγησης.»}
\ask{«Ποιο μικρό αλλά αντιπροσωπευτικό corpus θα ήταν αρκετό για τον πρώτο έλεγχο εφικτότητας των οικογενειών και των εξηγητικών ζευγών;»}
\newpage
% PAGE 5
\section*{4. M1 — γιατί υπάρχει το \(U\times F\)}

\subsection*{Η ερώτηση της M1}
Ποια representation διατηρεί την πληροφορία που χρειάζεται για να εντοπίζονται θέσεις μεταβλητού μήκους και να αναγνωρίζεται η οικογενειακή τους ταυτότητα σε νέα έργα/πηγές;

\subsection*{Οι τρεις μονάδες \(U\)}
\begin{tabularx}{\linewidth}{>{\bfseries}p{1.2cm}p{3.7cm}X}
\(u_1\) & Γραφικό σύμβολο & Μέγιστη πιστότητα στη γραφή.\\
\(u_2\) & Atomic melodic event & Κοινή λεπτομερής μελωδική κλίμακα και evaluation grid.\\
\(u_3\) & Deterministic compound neumatic unit & Ελέγχει αν μουσικολογικά/γραφικά σύνθετη μονάδα είναι αποδοτικότερη από την ατομική ανάλυση.\\
\end{tabularx}

Ο κανόνας \(u_3=f_{\mathrm{group}}(N)\) είναι σταθερός σε όλα τα επίπεδα \(F\), ώστε να αλλάζουν μόνο τα features και όχι η ίδια η ομαδοποίηση.

\subsection*{Η σωρευτική κλίμακα \(F\)}
\[
F=\{N,\ N+T,\ N+T+M,\ N+T+M+L\}.
\]
\textbf{N}: σημειογραφική/μελωδική πληροφορία. \textbf{T}: χρονική. \textbf{M}: τροπική/state. \textbf{L}: κειμενική/προσωδιακή.

Δεν πρόκειται για πλήρες \(2^3\) factorial των T/M/L. Είναι \textbf{διασταυρωμένη διάταξη 3×4 με διατεταγμένη σωρευτική κλίμακα γνωρισμάτων}. Ελέγχει την πρόσθετη αξία της πληροφορίας στη συγκεκριμένη σειρά.

\subsection*{Γιατί 12 και όχι 24 κύριες conditions}
Ο κύριος στόχος είναι selection, όχι πλήρης αιτιώδης αποσύνθεση κάθε interaction. Οι 12 περιορίζουν multiplicity και computational cost. Δύο post-selection diagnostics, \(N+M+L\) και \(N+T+L\), μαζί με το υπάρχον \(N+T+M\), ελέγχουν leave-one-group-out γύρω από το πλήρες \(F\) χωρίς να ξαναδιαλέγουν \(R^\star\).

\say{«Δεν λέω ότι τα T, M και L έχουν ανεξάρτητες main effects. Ελέγχω μια προκαθορισμένη σωρευτική κλίμακα πληροφορίας και κρατώ targeted diagnostics για το αν το πλήρες representation εξαρτάται ειδικά από T, M ή L.»}

\ask{«Θα θεωρούσατε αρκετό το 3×4 search space με targeted ablations, ή θα θέλατε ένα μικρό fractional factorial validation πριν κλειδώσει το τελικό paper της M1;»}

\newpage
% PAGE 6
\section*{5. M1 — fairness, \(R^\star\) και τι σημαίνει “καλύτερη”}

\subsection*{Fair comparison}
Οι representations έχουν διαφορετικό vocab και sequence length. Για αυτό η proposal χρησιμοποιεί κοινό atomic-event target για την information task και κοινό atomic-event grid για segmentation/evaluation. Στην boundary task χρειάζεται κοινό receptive field σε atomic events, όχι μηχανική μεταφορά του causal horizon \(H\) της next-event prediction.

\subsection*{Δύο μέτρα, δύο ρόλοι}
\begin{itemize}
\item \(M(R)\): primary boundary performance — micro boundary F1 στο κοινό πλήρως ελεγμένο development subset, με προκαθορισμένη ανοχή και one-to-one matching.
\item \(G(R)\): family performance — macro family recall σε externally documented occurrences, ώστε οι οικογένειες να έχουν ίσο βάρος.
\end{itemize}
Το work-balanced boundary F1 αναφέρεται ως diagnostic, όχι selector.

\subsection*{Το family readiness gate}
\[
g_{\min}=\max\{g_{\mathrm{use}},\widehat G_0^{\mathrm{pilot}}+\delta_{\mathrm{use}}\}.
\]
Οι \(g_{\mathrm{use}}\), \(\delta_{\mathrm{use}}\) δηλώνονται γραπτά και χρονοσημαίνονται πριν από την κύρια M1. Το \(g_{\min}\) δεν προσαρμόζεται αφού φανούν οι 12 συνθήκες.

\subsection*{Λειτουργική επιλογή}
\[
\widehat{\mathcal A}_{\mathrm{dec}}=\{R:\widehat G(R)\ge g_{\min}\},
\qquad
R^\star=\arg\max_{R\in\widehat{\mathcal A}_{\mathrm{dec}}}\widehat M(R).
\]
Η \(R^\star\) είναι λειτουργική απόφαση για τη M2. Δεν αποτελεί απόδειξη μοναδικής πληθυσμιακής υπεροχής. Το \(\mathcal R^\star\) αναφέρει selection uncertainty. Η M2 δεν ξαναδιαλέγει representation.

\subsection*{Τι κόστος έχει η προτεραιότητα στην οριοθέτηση;}
\[
\Delta_G^{\mathrm{loss}}=G(R^\star)-\max_R G(R)\le0.
\]
Αναφέρεται υποχρεωτικά ώστε να φαίνεται πόση family performance θυσιάζει η boundary-optimal admissible representation.

\say{«Χωρίζω decision από inference. Χρειάζομαι μία operational representation για να προχωρήσει η M2, αλλά αναφέρω χωριστά την αβεβαιότητα και δεν παρουσιάζω την πρώτη θέση ως απόδειξη μοναδικής υπεροχής.»}

\ask{«Θα προτιμούσατε το primary boundary F1 να παραμείνει micro με work-balanced diagnostic ή να γίνει work-level primary λόγω του μικρού αριθμού ανεξάρτητων έργων;»}

\newpage
% PAGE 7
\section*{6. M2 — γιατί EISE και όχι απλός classifier}

\subsection*{Το μαθησιακό πρόβλημα}
Η M2 πρέπει ταυτόχρονα να χειριστεί:
\begin{itemize}
\item θέσεις μεταβλητού μήκους,
\item αβέβαια/μερικά όρια,
\item ελλιπείς οικογενειακές ετικέτες,
\item open-set περιπτώσεις,
\item κοινή οικογενειακή ταυτότητα,
\item διατήρηση της εσωτερικής πραγμάτωσης,
\item ελεγχόμενη σύγκριση δύο πηγών relational supervision.
\end{itemize}

\subsection*{Γιατί semi-Markov}
Ένα linear CRF δουλεύει κυρίως σε event-level labels. Το semi-Markov CRF βαθμολογεί ολόκληρα spans μεταβλητού μήκους και μπορεί να εκτιμήσει όρια και labels σε κοινό structured model. Η μηδενικής τάξης αρχική μορφή αποφεύγει να κερδίζει το family identification από transition priors οικογενειών.

\subsection*{Γιατί partial likelihood}
Η απουσία annotation δεν ισοδυναμεί με αρνητική ένδειξη. Για κάθε window επιτρέπεται σύνολο συμβατών structured outputs και η loss αθροίζει την πιθανότητα αυτών των επιτρεπτών αναλύσεων. Έτσι αμφίσημα όρια και μερικές ετικέτες δεν μετατρέπονται τεχνητά σε λάθη.

\subsection*{Γιατί open set}
Έγκυρο μουσικό τμήμα μπορεί να μην αντιστοιχεί με ασφάλεια σε γνωστή οικογένεια. Για αυτό υπάρχει \(\varnothing_F\) και novelty score. Η αποχή είναι ουσιώδης, ειδικά στη M5, γιατί αποτρέπει την αναγκαστική ένταξη κάθε νέο σχήματος σε γνωστή οικογένεια.

\say{«Η πολυπλοκότητα του EISE δεν ξεκινά από την επιθυμία για ένα σύνθετο neural model. Προκύπτει από τέσσερις ιδιότητες του προβλήματος: variable length, partial labels, open set και ανάγκη να διαχωρίζεται identity από realization.»}

\ask{«Στην πρώτη implementation θα επιμένατε σε zero-order semi-CRF ή θα θεωρούσατε ασφαλέστερο να ξεκινήσει η M2 με απλούστερο CRF+segment pooling και το semi-CRF να κερδίσει τη θέση του σε περιορισμένο development comparison;»}

\newpage
% PAGE 8
\section*{7. EISE — αρχιτεκτονική που πρέπει να μπορείς να σχεδιάσεις στον πίνακα}

\subsection*{Pipeline}
\[
X\xrightarrow{R^\star} h_1,\ldots,h_T
\xrightarrow{\text{candidate spans}} r_s
\xrightarrow{q_{id}} z_{id}(s)
\xrightarrow{\text{prototype/open-set}} \phi_\theta(s,F).
\]

Για span \(s=(a,b)\):
\[
r_s=\left[\frac{1}{b-a+1}\sum_{t=a}^{b}h_t;\ h_a;\ h_b\right],
\qquad
z_{\mathrm{id}}(s)=\operatorname{Normalize}(q_{\mathrm{id}}(r_s)).
\]

Για γνωστή οικογένεια:
\[
\phi_\theta(s,F)=\psi_{\mathrm{seg}}(r_s)-\frac{d(z_{\mathrm{id}}(s),p_F)}{\tau}.
\]
Για έγκυρο αλλά άγνωστο τμήμα υπάρχει ειδική \(\varnothing_F\) με novelty function.

\subsection*{Τέσσερις όροι της loss}
\[
\mathcal L_{\mathrm{EISE}}=
\mathcal L_{X,\mathrm{seg}}^{\mathrm{partial}}
+\lambda_{\mathrm{rel}}\mathcal L_{\mathrm{pairs}}^{\mathrm{id}}
+\lambda_{\mathrm{proto}}\mathcal L_{\mathrm{proto}}
+\lambda_{\mathrm{loc}}\mathcal L_{\mathrm{local}}.
\]
\begin{itemize}
\item \(L_{seg}^{partial}\): boundaries/labels με partial supervision.
\item \(L_{pairs}^{id}\): φέρνει relational positives κοντά στο identity space.
\item \(L_{proto}\): οργανώνει τις γνωστές οικογένειες γύρω από prototypes.
\item \(L_{local}\): διατηρεί πληροφορία συγκεκριμένης πραγμάτωσης στα local states.
\end{itemize}

\subsection*{Κρίσιμη διάκριση}
Το \(z_{id}\) είναι η συμπυκνωμένη οικογενειακή ταυτότητα. Τα \(h_t\) διατηρούν τοπική πραγμάτωση. Για αυτό η M5 χρησιμοποιεί local states για \(D_{within}\) και όχι μόνο \(z_{id}\).

\subsection*{Τι δεν έχει αποδειχθεί ακόμη}
Η τρέχουσα proposal επιλέγει τον EISE ως βασική αρχιτεκτονική σχεδιασμού. Η τελική μορφή structured output, hyperparameters και ορισμένα diagnostics οριστικοποιούνται μετά από pilot/dev tests. Δεν υπάρχουν ακόμη experimental results που να δικαιολογούν φράσεις τύπου «ο EISE υπερέχει».

\say{«Το \(z_{id}\) πρέπει να ξεχνά μέρος της ενδοοικογενειακής διαφοράς. Τα \(h_t\) πρέπει να κρατούν αρκετή local πληροφορία ώστε να μην καταρρεύσει η realization geometry. Αυτή η διπλή απαίτηση είναι ο λόγος που έχω identity projection και local objective.»}

\newpage
% PAGE 9
\section*{8. B0, B1, B2, EISE — η καρδιά του ελεγχόμενου πειράματος}

\begin{tabularx}{\linewidth}{>{\bfseries}p{1.25cm}p{4.2cm}X}
B0 & Άμεσο συμβολικό ιστορικό profile & Τι εξηγείται χωρίς representation learning.\\
B1 & Ίδια neural/segmental υποδομή χωρίς relational supervision & Ελέγχει τι δίνει η κοινή αρχιτεκτονική και η οικογενειακή εποπτεία.\\
B2 & B1 + matched within-family positive pairs & Ισχυρός μη ιστορικός relational comparator.\\
EISE & Ίδια υποδομή + documented explicit-exegesis positive pairs & Ελέγχει την πρόσθετη αξία της προέλευσης της σχέσης.\\
\end{tabularx}

\subsection*{Κύριες αντιθέσεις}
\[
\Delta_{\mathrm{hist}}=S_{\mathrm{EISE}}-S_{B2},
\qquad
\Delta_{\mathrm{learn}}=S_{\mathrm{EISE}}-S_{B0}.
\]
Το \(S_m\) είναι συμφωνία του ιστορικού profile του μοντέλου με ανεξάρτητο performance profile.

\subsection*{Τι απομονώνει το \(\Delta_{hist}\)}
Μόνο εφόσον B2 και EISE είναι πραγματικά matched σε architecture, data, exposure, pair count/difficulty, compute, preprocessing και family information, το κύριο ερώτημα γίνεται:

\begin{center}
\textbf{Προσθέτει η ιστορική προέλευση των positive relations κάτι πέρα από matched family relation supervision;}
\end{center}

\subsection*{Σενάρια αποτελέσματος}
\begin{itemize}
\item \(\Delta_{hist}>0\): οι documented exegesis relations προσθέτουν χρήσιμη οργάνωση υπό το συγκεκριμένο external criterion.
\item \(\Delta_{hist}\approx0\): δεν τεκμηριώνεται πρόσθετη αξία πέρα από matched family supervision. Δεν σημαίνει ότι «η μορφή αρκεί» ούτε ότι οι εξηγήσεις είναι μουσικολογικά άχρηστες.
\item \(\Delta_{hist}<0\): η εξηγητική relation source, όπως επιχειρησιακά ορίστηκε, δεν βελτιώνει το criterion και μπορεί να είναι αραιή, ετερογενής ή να κωδικοποιεί άλλη σχέση.
\end{itemize}

\avoid{«Η B2 δεν μπορεί να μάθει εσωτερική σταθερότητα ενώ ο EISE μπορεί.» Η B2 μπορεί να μάθει regularities από τη distribution των same-family pairs. Η διαφορά είναι η provenance της relational supervision.}

\ask{«Θα θέλατε το \(B2_{sim}\) ως υποχρεωτικό stricter comparator ή μόνο εφόσον η απογραφή δείξει επαρκή coverage ώστε να μην αποδυναμωθεί η ισχύς της κύριας B2--EISE σύγκρισης;»}

\newpage
% PAGE 10
\section*{9. M3a — ανεξάρτητη μουσικολογική επικύρωση}
\subsection*{Διαφορετικός στόχος από το ΠΕ0}
Το ΠΕ0 κατασκευάζει και ελέγχει το λεξιλόγιο \emph{ανάπτυξης}. Η M3a ελέγχει αν άλλοι γνώστες της παράδοσης εφαρμόζουν το ίδιο μουσικολογικό πλαίσιο σε δεδομένα που δεν χρησίμευσαν για βελτιστοποίηση. Η M3a μπορεί να προχωρά παράλληλα με την M2, διατηρώντας ανεξάρτητα δεδομένα.
\subsection*{Η τυφλή μελέτη}
\begin{itemize}
\item 3--5 ανεξάρτητοι ειδικοί με επιδίωξη ποικιλίας ψαλτικής/ερευνητικής προέλευσης.
\item Βλέπουν το αναγκαίο μουσικό συμφραζόμενο και τις οριστικοποιημένες οδηγίες, \emph{όχι} τις προβλέψεις του EISE ή τις αποφάσεις των άλλων.
\item Επισημειώνουν όρια, οικογενειακή ταυτότητα και, όπου τεκμηριώνεται, λειτουργικό ρόλο· έχουν επιλογή \emph{ανεπίλυτο}.
\item Μετράμε χωριστά συμφωνία για όρια και για οικογενειακή απόδοση, μαζί με την έκταση των ανεπίλυτων περιπτώσεων.
\item Δεν αλλάζουμε αναδρομικά την προδιαγραφή του μοντέλου ή το test set ώστε να αυξηθεί η συμφωνία.
\end{itemize}
\subsection*{Τι σημαίνει καλό αποτέλεσμα;}
Ο συνδυασμός συμφωνίας μεταξύ ειδικών, ανεξάρτητης κριτικής των πηγών και καλής ανάκτησης σε νέα έργα στηρίζει τη χρησιμότητα του λεξιλογίου. Δεν τεκμηριώνει ότι κάθε οικογένεια φέρει μία ενιαία ιστορική ονομασία. Η διαφωνία μπορεί να δείχνει πραγματική μεταβλητότητα της παράδοσης.
\say{«Δεν ζητώ από τους ειδικούς να επικυρώσουν τις ετικέτες του EISE. Ζητώ ανεξάρτητες κρίσεις για τα ίδια μουσικά χωρία και ελέγχω τη σχέση τους με τα υπολογιστικά αποτελέσματα.»}
\avoid{«Τυφλή σημαίνει ότι οι ειδικοί δεν ξέρουν τίποτε για το μέλος ή τον ήχο» ή «η M3a ξαναγράφει τις ετικέτες της M2».}
\newpage
% PAGE 11
\section*{10. M3c — μουσική προσδοκία, PPM και ανθρώπινη κρίση}
\subsection*{Τι εξετάζω}
Ο ακροατής ακούει ένα μουσικό απόσπασμα και τη συνέχεια και βαθμολογεί πόσο απρόσμενη του φαίνεται. Το ζητούμενο είναι αν η οικογένεια της συνέχειας έχει πρόσθετη ερμηνευτική αξία σε σχέση με τα μεμονωμένα γεγονότα.
\subsection*{Δύο διαφορετικά πληροφοριακά μεγέθη}
\[
I_F=-\log_2 P(F_t\mid F_{<t},M,L),\qquad I_{event}=\text{σύνοψη της έκπληξης των γεγονότων}.
\]
Το \(I_F\) αφορά πόσο απρόσμενη είναι η οικογένεια \emph{μετά την προηγούμενη διαδοχή} και με δεδομένο τροπικό/δομικό/κειμενικό περιβάλλον. Το \(I_{event}\) αφορά τα μεμονωμένα γεγονότα της ίδιας συνέχειας, με κύρια σύνοψη τον μέσο surprisal τους. Δεν συγκρίνουμε απλώς τις ακατέργαστες αριθμητικές τιμές τους.
\subsection*{Γιατί PPM και ποιος ο υπολογιστικός έλεγχος;}
Το PPM χρησιμοποιεί προηγούμενο συμφραζόμενο μεταβλητού μήκους και, όταν αυτό είναι σπάνιο, επιστρέφει σε μικρότερο. Ελέγχεται σε \emph{έργα εκτός εκπαίδευσης} έναντι του μοντέλου \(P(F_t\mid M,L)\), ώστε να φανεί αν \emph{η ακολουθία} προσφέρει πληροφορία πέρα από τον ήχο/δομικό τόπο.
\subsection*{Γιατί υψηλό/χαμηλό στους αντίθετους συνδυασμούς;}
Χρειαζόμαστε περιπτώσεις \(I_F\uparrow,I_{event}\downarrow\) και \(I_F\downarrow,I_{event}\uparrow\), όχι για να «πάμε στα άκρα», αλλά για να χωρίσουμε τις επιδράσεις. Η μουσική φυσικότητα ελέγχεται τυφλά. Αν οι περιπτώσεις δεν επαρκούν, η κύρια υπόθεση δεν εξαναγκάζεται.
\subsection*{Εμπειρία και λατρευτικό πλαίσιο}
Ο \(E_{Byz}\) είναι διαβαθμισμένος δείκτης ψαλτικής εξοικείωσης, όχι απλώς ψάλτης/μη ψάλτης. Το \(R\) είναι πειραματική \emph{μη μουσική ένδειξη} λατρευτικού πλαισίου, όχι αναπαράσταση ολόκληρης της επιτέλεσης. Το \(\Delta_{ctx}\) εξετάζει αν αλλάζει διαφορετικά η σχέση των κρίσεων με \(I_F\) και \(I_{event}\).
\say{«Η M3c ελέγχει σύγχρονη ανθρώπινη μουσική προσδοκία. Αρνητικό αποτέλεσμα δεν αναιρεί ότι οι θέσεις είναι χρήσιμες μουσικολογικά/υπολογιστικά.»}
\newpage
% PAGE 12
\section*{11. M3b — κατηγοριοποίηση πέρα από επιφανειακή ομοιότητα}
\subsection*{Τρία αποσπάσματα, ένα ερώτημα}
Ο ακροατής ακούει τρεις πραγματώσεις και επιλέγει ποια ταιριάζει λιγότερο με τις άλλες δύο. Σε μεικτή τριάδα έχουμε δύο πραγματώσεις της ίδιας οικογένειας και μία άλλης. Περιλαμβάνονται επίσης τριάδες με τρεις πραγματώσεις της ίδιας οικογένειας, για τη μελέτη της εσωτερικής διαφοροποίησης.
\subsection*{Το ερευνητικό πρόβλημα}
Αν η οικογενειακή συγγένεια συμπίπτει πάντα με τη μεγαλύτερη ακουστική ομοιότητα, μια σωστή απάντηση δεν αποδεικνύει κατηγοριοποίηση. Γι' αυτό επιλέγουμε τριάδες από \emph{κοινή περιοχή συμβολικών και ακουστικών αποστάσεων} και περιλαμβάνουμε δύσκολα παραδείγματα: ίδια οικογένεια με μεγάλη επιφανειακή διαφορά, διαφορετικές οικογένειες με συγκρίσιμη επιφανειακή ομοιότητα.
\subsection*{Τι δείχνει η H\textsubscript{cat};}
Εξετάζουμε αν η οικογενειακή ένταξη προβλέπει την ανθρώπινη επιλογή \emph{αφού ελεγχθούν} οι επιφανειακές αποστάσεις και άλλες προσυμφωνημένες μεταβλητές. Μια θετική σχέση δεν αποδεικνύει ότι ο ακροατής γνωρίζει ιστορικές ονομασίες ή εξηγητικές σχέσεις.
\subsection*{Σημερινό καθεστώς και πιθανή απλοποίηση}
\textbf{Στη v8.71a η M3b παραμένει επιβεβαιωτική υπόθεση} με προκαθορισμένα κριτήρια εφικτότητας, και ακολουθεί τη M3c στην ανθρώπινη συλλογή. Η πιθανή μεταφορά της σε προαιρετική βαθμίδα είναι \emph{συζήτηση για μελλοντική αναθεώρηση}, όχι ήδη εγκεκριμένη αλλαγή. Αν γίνει, αλλάζουν αντίστοιχα προεγγραφή, στατιστική ισχύς και οικογένεια ελέγχων.
\say{«Η M3a ελέγχει αν το λεξιλόγιο είναι μουσικολογικά αναγνωρίσιμο· η M3b αν η οικογενειακή διάκριση έχει αντιληπτή κατηγορική οργάνωση πέρα από την απλή ομοιότητα.»}
\ask{«Θα κρατούσατε την πλήρη M3b ως επιβεβαιωτική δοκιμή ή θα τη μεταφέρατε μετά τον πιλότο σε επιλέξιμη ενίσχυση, χωρίς να αποδυναμωθεί η M3c;»}
\newpage
% PAGE 13
\section*{12. M4 — γραφή, εκτέλεση και ανεξάρτητο κριτήριο}
\subsection*{Πρώτα αντιστοίχιση, έπειτα συμπέρασμα}
Χρειαζόμαστε \textbf{πολλαπλές ηχογραφήσεις του ίδιου συγκρίσιμου μουσικού υλικού}. Ευθυγραμμίζουμε χειροκίνητα τα όρια συλλαβών και, όπου η ακρίβεια το επιτρέπει, αντιστοιχίζουμε εντός της συλλαβής τα επιμέρους μελωδικά γεγονότα με score-informed DTW ή μοντέλο καταστάσεων. Μια μελισματική συλλαβή δεν είναι ένα γεγονός. Αν η λεπτομερής ευθυγράμμιση αποτύχει, εργαζόμαστε σε συλλαβή/νευματική ομάδα.
\subsection*{Η κρίσιμη διάκριση}
Μια πρόσθετη μελωδική ανάπτυξη δεν πρέπει να απορροφηθεί ως «καθυστέρηση» σε έναν απλό χρονικό μετασχηματισμό. Την επισημαίνουμε ως δομική/μελωδική διαφοροποίηση όταν τα διαθέσιμα δεδομένα την τεκμηριώνουν. Χωριστά καταγράφουμε κανονικοποιημένο χρονισμό, συνολική αγωγή, τοπική αγωγή, καταληκτική επιμήκυνση και, όπου αξιόπιστο, τονικό ύψος.
\subsection*{Η κύρια εξωτερική σύγκριση της M2}
\begin{itemize}
\item \(A(z)\): ιστορικό προφίλ παραλλαγής ανά εσωτερική περιοχή \(z\) της θέσης.
\item \(\Gamma(z)\): ανεξάρτητη διασπορά μεταξύ σύγχρονων εκτελεστών στην ίδια μουσική περιοχή.
\item \(S_m\): συμφωνία κατάταξης ανάμεσα σε ιστορικό και εκτελεστικό προφίλ για \(m\in\{B0,B2,EISE\}\).
\end{itemize}
Η προεπιλεγμένη κύρια εκτελεστική μεταβλητή είναι η διασπορά \emph{κανονικοποιημένου λογαριθμικού χρονισμού}, αν η ποιότητά της επαρκεί. Το ύψος είναι δευτερεύον. Τα εκτελεστικά δεδομένα δεν αναπροσαρμόζουν εκ των υστέρων τον EISE.
\subsection*{Όριο ιστορικού ισχυρισμού}
Η σύμπτωση μιας σύγχρονης εκτελεστικής μορφής με μία ιστορική εξήγηση είναι μορφολογική/εκτελεστική αντιστοιχία· \emph{δεν} αποδεικνύει συγκεκριμένη διδασκαλική γενεαλογία ή ανακατασκευή του ιστορικού ακούσματος.
\say{«Η M4 είναι εξωτερικός έλεγχος της οργανωμένης παραλλαγής, όχι ακουστική ανακατασκευή του 19ου αιώνα.»}
\newpage
% PAGE 14
\section*{13. M5 — από τη θέση στο πλήρες έργο}

\subsection*{Η σωστή εννοιολογική εικόνα}
Η θέση είναι η \textbf{κεντρική ενδιάμεση μονάδα}, όχι η μοναδική μονάδα της πλήρους αναπαράστασης. Atomic events και συλλαβές βρίσκονται χαμηλότερα. Ο χάρτης εμφανίσεων, η διαδοχή, η λειτουργική τοποθέτηση και το πλήρες έργο βρίσκονται υψηλότερα. Το έργο δεν αντιμετωπίζεται ως bag of positions ούτε ως ένα αδιαφοροποίητο embedding.

\[
\begin{aligned}
\text{events/syllables} &\rightarrow \text{position occurrence}\\
&\rightarrow \{\text{family identity, local realization}\}\\
&\rightarrow \text{sequence/function}\rightarrow \text{whole-work profile}.
\end{aligned}
\]

\subsection*{Οι τέσσερις διαστάσεις}
\[
\mathbf D(w;\Omega)=
\bigl(D_{fam},D_{seq},D_{func},D_{within}\bigr).
\]
\begin{itemize}
\item \(D_{fam}\): χρησιμοποιούνται οι ίδιες οικογένειες και με παρόμοια κατανομή;
\item \(D_{seq}\): πώς διατάσσονται και πόσο προβλέψιμες είναι οι ακολουθίες;
\item \(D_{func}\): χρησιμοποιούνται οι οικογένειες στους ίδιους δομικούς/κειμενικούς/τροπικούς τόπους;
\item \(D_{within}\): όταν χρησιμοποιείται η ίδια οικογένεια, πώς διαφέρει η συγκεκριμένη πραγμάτωση;
\end{itemize}

\subsection*{Γιατί δεν συγχωνεύονται}
Έργο μπορεί να έχει συμβατικό λεξιλόγιο αλλά ασυνήθιστη σειρά, ή συμβατική σειρά αλλά ιδιότυπη πραγματοποίηση. Η M5 επομένως παράγει πολυδιάστατο profile, όχι «ποσοστό κλασικότητας».

\subsection*{Τοπική ερμηνευσιμότητα}
Το \(D_{within}\) παραμένει ενιαίο. Όπου υπάρχουν αξιόπιστες κοινές εσωτερικές περιοχές, συνοδεύεται από \(Q_x,p_x,d_{time}\) και τοπικά profiles. Δεν επιβάλλεται καθολικό σχήμα core/extension/cadence.

\subsection*{Η πραγματική δύναμη της M5}
\textbf{Global profile + local traceability.} Κάθε διαφορά πρέπει να μπορεί να επιστρέψει σε οικογένεια, εμφάνιση, έργο, source locus και συγκεκριμένο σημείο της πραγμάτωσης.

\say{«Η M5 δεν κρίνει αν ένα έργο είναι “καλό”. Χαρακτηρίζει πού και με ποιον τρόπο διαφοροποιείται από ένα ρητά ορισμένο ιστορικό περιβάλλον αναφοράς.»}

\newpage
% PAGE 15
\section*{14. M5 — τι ακριβώς μπορείς και δεν μπορείς να ισχυριστείς}

\subsection*{Το περιβάλλον αναφοράς \(\Omega\)}
\[
\Omega=(D_\Omega,\rho,\Pi,\mathcal K_\Omega).
\]
Περιλαμβάνει ρητό corpus αναφοράς, οριστικοποιημένη representation/family assignment, sequence model και κανόνες ένταξης/χρονολόγησης/ανεξαρτησίας. Άρα κάθε claim είναι σχετικό με συγκεκριμένο \(\Omega\), όχι με αφηρημένη «κλασική παράδοση».

\subsection*{M5-min και M5-full}
Η M5-min δεσμεύεται για \(D_{fam}\) και \(D_{within}\) και για χάρτη πλήρους έργου. Η M5-full προσθέτει \(D_{seq}\) και \(D_{func}\) όταν υπάρχει επαρκής coverage ακολουθιών, λειτουργικών τόπων και ανεξάρτητων μελοποιών.

\subsection*{Για Νικόλαο}
Η ισχυρή ερώτηση δεν είναι «πόσο αποκλίνει από τις κλασικές θέσεις;» αλλά:
\begin{quote}
Σε ποιες διαστάσεις και σε ποια συγκεκριμένα χωρία διαφοροποιείται ο Νικόλαος από το δηλωμένο σώμα αναφοράς: στο λεξιλόγιο οικογενειών, στη διαδοχή, στη λειτουργική χρήση ή στην εσωτερική πραγμάτωση;
\end{quote}
Αυτό επιτρέπει να ξεχωρίσεις διαφορετικό λεξιλόγιο από διαφορετική οργάνωση του ίδιου λεξιλογίου και από διαφορετική ανάπτυξη των ίδιων οικογενειών.

\subsection*{Τι δεν βάζεις στη M5}
\begin{itemize}
\item Ενιαίο score ποιότητας/ορθότητας/αισθητικής αξίας.
\item Αυτόματη ταύτιση απόκλισης με λάθος ή «μη κλασικό».
\item Kama muta ως core M5 dependent variable. Η σχέση με ανθρώπινη εμπειρία ανήκει σε M3c ή exploratory bridge.
\item Υποχρεωτική φράση/περίοδος ως νέο latent layer χωρίς νέο annotation/model.
\end{itemize}

\ask{«Για την πρώτη δημοσίευση της M5, θα προτιμούσατε να μείνω αυστηρά στο M5-min με πλήρη traceability ή να επενδύσω από την αρχή σε sequence/function, με κίνδυνο αραιότερης υποστήριξης;»}

\newpage
% PAGE 16
\section*{15. Στατιστική πειθαρχία, leakage και τι πρέπει να ξέρεις χωρίς να γίνεις στατιστικός}

\subsection*{Τρεις αρχές}
\begin{enumerate}
\item \textbf{Selection μόνο σε development.} Test/historical challenge/human outcomes δεν αποφασίζουν representation ή hyperparameters.
\item \textbf{Grouped splits.} Ολόκληρα έργα και συγγενείς πηγές μένουν μαζί ώστε παραλλαγές του ίδιου υλικού να μη διαρρέουν μεταξύ train/test.
\item \textbf{Matched contrasts.} Στη B2--EISE κάθε κύρια διαφορά πρέπει να είναι η source/provenance των positive relations, όχι compute, exposure ή διαφορετική αρχιτεκτονική.
\end{enumerate}

\subsection*{Αβεβαιότητα}
Η proposal χρησιμοποιεί cluster/group-aware resampling, matched training seeds, simultaneous intervals όπου υπάρχει οικογένεια συγκρίσεων και αποφεύγει να μεταφράζει non-significance σε equivalence.

\subsection*{Τι πρέπει να μπορείς να πεις για power}
Δεν υπόσχεσαι αριθμούς πριν από την απογραφή. Το ΠΕ0 εκτιμά cluster counts, coverage, pair counts, variance across seeds και reliability. Όπου δεν υπάρχει estimability, το αποτέλεσμα υποβαθμίζεται σε descriptive/exploratory αντί να χαλαρώνει post hoc το criterion.

\subsection*{Αν \(\widehat{\mathcal A}_{dec}=\varnothing\)}
Δεν χαμηλώνεις το \(g_{min}\) μετά τα αποτελέσματα και δεν ξεκινά confirmatory M2. Μπορεί να χρησιμοποιηθεί diagnostic best-boundary representation για pipeline/error analysis, χωρίς να υποκαταστήσει \(R^\star\). Η M3a μπορεί να συνεχίσει ανεξάρτητα· οι M3b/M3c μόνο με τα δικά τους gates· η κανονική M5 απαιτεί \(R^\star\).

\say{«Οι πύλες δεν υπάρχουν για να σώσουν θετικά αποτελέσματα. Υπάρχουν για να ορίσουν πριν από την ανάλυση πότε ένα estimand είναι πραγματικά εκτιμήσιμο.»}

\ask{«Ποια από τα statistical safeguards θεωρείτε ουσιώδη για proposal stage και ποια θα προτιμούσατε να μείνουν ως analysis-plan detail μετά το ΠΕ0 ώστε να μη φαίνεται υπερσχεδιασμένη η πρόταση;»}

\newpage
% PAGE 17
\section*{16. Διάταξη διατριβής, εξαρτήσεις και 48 μήνες}
\subsection*{Οι βαθμίδες είναι σκόπιμα διαφορετικές}
\begin{itemize}
\item \textbf{Αναγκαίος υπολογιστικός/μουσικολογικός πυρήνας}: M1, M2, M3a και το ελάχιστο ανεξάρτητο εκτελεστικό κριτήριο της M4 για την M2.
\item \textbf{Ανθρώπινη επιβεβαιωτική βαθμίδα στη v8.71a}: M3c, έπειτα M3b, με χωριστούς ελέγχους επάρκειας και κοινό έλεγχο πολλαπλότητας.
\item \textbf{Προγραμματισμένο ιστορικό παραδοτέο}: M5-min με \(D_{fam}\), \(D_{within}\), χάρτη εμφανίσεων, αβεβαιότητα και συγκεκριμένα χωρία· πλήρης \(D_{seq}\), \(D_{func}\) αν το επιτρέπουν τα δεδομένα.
\item \textbf{Επιπλέον}: ευρύτερη M4, λεπτομερέστερες ιστορικές ομάδες, δευτερεύουσες αναλύσεις και M3ν (EEG/MEG) μόνο υπό σαφείς προϋποθέσεις.
\end{itemize}
\subsection*{Χρονική ακολουθία, όχι «όλα μετά από όλα»}
ΠΕ0/ΠΕ1: απογραφή, άδειες, σώμα, συμβολικό schema, κάρτες και πιλοτική αξιοπιστία. M1: κατάλληλη αναπαράσταση και \(R^\star\). M3a: χωριστή επικύρωση, παράλληλα με την ανάπτυξη. M2: εκπαιδευτική σύγκριση B2--EISE και, εφόσον πληρούνται οι προϋποθέσεις, κύρια \(\Delta_{hist}\). M3c πριν από M3b στην ανθρώπινη συλλογή. M4: αρχικά ελάχιστο ανεξάρτητο κριτήριο, κατόπιν ευρύτερη ανάλυση. M5: ιστορική εφαρμογή Νικολάου. Σύνθεση έως 48 μήνες.
\subsection*{Τι έχει οριστικοποιηθεί ως σχέδιο και τι όχι}
\textbf{Ορισμοί/πρωτόκολλο}: διάκριση οικογένειας--εξήγησης, \(U\times F\) 3×4, κανόνας \(R^\star\), αντιστοιχισμένη B2--EISE, ανεξάρτητα κριτήρια, M5-min. \textbf{Ανοιχτά}: πραγματική κάλυψη, μέγεθος corpus, αξιοπιστία alignment, εφικτότητα ανθρώπινων ερεθισμάτων, εφαρμοστέα κλίμακα, αριθμός επαναλήψεων και τελικές υπερπαράμετροι.
\subsection*{Ιδιαίτερο στοιχείο επιστημονικής συζήτησης}
Ο δυνητικός επιβλέπων έχει λάβει παλαιότερο κείμενο. Παρουσιάζω πρώτα τη σταθερή βασική ιδέα και περιγράφω τη v8.71a ως \emph{μεταγενέστερη αποσαφήνιση πρωτοκόλλου}, όχι ως τετελεσμένα πειραματικά αποτελέσματα.
\say{«Το σχέδιο είναι συγκεκριμένο ως προς τις επιστημονικές αντιθέσεις, αλλά η πρώτη χρονιά πρέπει να ελέγξει πού τα πραγματικά δεδομένα επιτρέπουν αξιόπιστη υλοποίηση.»}
\newpage
% PAGE 18
\section*{17. Πιθανές δύσκολες ερωτήσεις του επιβλέποντα — σύντομες απαντήσεις}

\small
\begin{tabularx}{\linewidth}{>{\bfseries}p{5.0cm}X}
Γιατί θέση και όχι event Transformer; & Ο event Transformer είναι ισχυρό baseline. Το ερώτημα είναι αν η explicit variable-length θέση προσθέτει boundaries, family identity και interpretability πέρα από event-level modelling. Αν όχι, αυτό είναι έγκυρο αποτέλεσμα.\\[2mm]
Γιατί όχι BPE/learned tokenization; & Τα χρησιμοποιώ ως data-driven comparator όπου είναι εφικτό. Η θέση όμως είναι musicological hypothesis, όχι απλώς compression unit.\\[2mm]
Γιατί semi-CRF; & Επειδή το label ανήκει σε variable-length span και τα boundaries είναι μέρος του inference. Το zero-order design περιορίζει shortcuts από family transitions.\\[2mm]
Γιατί EISE; & Χρειάζομαι κοινή architecture όπου B2/EISE διαφέρουν κυρίως στην provenance των positive relations.\\[2mm]
Τι αν \(\Delta_{hist}=0\); & Τότε δεν τεκμηριώνεται πρόσθετη αξία της historical relation source πέρα από matched family supervision στο συγκεκριμένο criterion. Το PhD εξακολουθεί να δίνει M1, M3a, human/performance findings και M5.\\[2mm]
Τι αν event model είναι καλύτερο; & Τότε η explicit θέση δεν κερδίζει σε εκείνο το task. Εξετάζω αν διατηρεί άλλες αξίες — family validity, retrieval, interpretability — χωρίς να βαφτίζω αποτυχία ως επιτυχία.\\[2mm]
Τι αν οι οικογένειες είναι ασαφείς; & Partial annotations, multiple acceptable analyses, open-set/abstention και independent M3a. Δεν μετατρέπω αμφισημία σε hard negatives.\\[2mm]
Γιατί σύγχρονες εκτελέσεις για ιστορικό question; & Ως external criterion της οργανωμένης variation, όχι ως reconstruction ιστορικής εκτέλεσης.\\[2mm]
Γιατί τόσο σύνθετο PhD; & Η ιεραρχία έχει core/minimal/full/conditional επίπεδα. Ο βασικός πυρήνας ολοκληρώνεται χωρίς να ενεργοποιηθούν όλες οι επεκτάσεις.\\[2mm]
Πόσο είναι δικό σου το computational design; & Η αρχιτεκτονική διαμορφώθηκε επαναληπτικά με βιβλιογραφία, κριτική και computational assistance. Η ευθύνη είναι να μπορώ να αιτιολογήσω, να υλοποιήσω και να ελέγξω κάθε τελική επιλογή.\\
\end{tabularx}
\normalsize

\begin{redline}
\textbf{Κόκκινη γραμμή:} Μην απαντήσεις με ονόματα μοντέλων χωρίς να λες ποιο methodological problem λύνουν. Η καλύτερη φόρμα είναι: «Το πρόβλημα είναι Χ. Η μέθοδος Υ το χειρίζεται επειδή Ζ. Η επιλογή δεν θεωρείται δεδομένη· ελέγχεται έναντι Β με criterion C.»
\end{redline}

\newpage
% PAGE 19
\section*{18. Αντίλογος που πρέπει να έχεις έτοιμο}

\subsection*{«Η θέση είναι πολύ ασαφής για computation.»}
Απάντηση: ακριβώς για αυτό το PhD δεν την παίρνει ως δεδομένη. Ορίζει candidate spans, partial references, multiple acceptable boundaries, family validation και data-driven comparators. Η αμφισημία γίνεται μέρος του μοντέλου/αξιολόγησης.

\subsection*{«Το \(U\times F\) είναι αυθαίρετο.»}
Απάντηση: οι τρεις unit scales ελέγχουν surface/atomic/compound encoding. Η feature ladder είναι προκαθορισμένη και σωρευτική. Το πλήρες factorial δεν είναι ο στόχος, αλλά υπάρχουν targeted post-selection ablations. Η επιλογή γίνεται σε κοινό task/grid.

\subsection*{«Το EISE είναι overengineered.»}
Απάντηση: variable spans + partial labels + open set + identity/realization + matched relational contrast απαιτούν περισσότερα από έναν απλό classifier. Παρ' όλα αυτά, η αρχιτεκτονική έχει simpler comparators και pilot gates.

\subsection*{«Δεν μπορείς να αποδώσεις causal αξία στην εξήγηση.»}
Απάντηση: σωστά. Το \(\Delta_{hist}\) είναι controlled predictive/representational contrast της provenance των positive relations υπό matching, όχι causal claim για ιστορική διαδικασία.

\subsection*{«Η M5 κινδυνεύει να γίνει αξιολόγηση ποιότητας.»}
Απάντηση: η M5 παράγει reference-relative, multidimensional differentiation profile. Απόκλιση δεν ισοδυναμεί με λάθος ή χαμηλή αισθητική αξία. Κάθε claim περιορίζεται στο \(\Omega\).

\subsection*{«Έχεις πάρα πολλά πράγματα.»}
Απάντηση: core = M1+M2+M3a+minimum M4. M3 human family και M5-min είναι επόμενες βαθμίδες. Full M4, full M5 και νευρική επέκταση ενεργοποιούνται μόνο με feasibility.

\subsection*{«Η ιστορική ακρόαση δεν είναι η σημερινή.»}
Απάντηση: η πρόταση διαχωρίζει historical sources, computational structure και contemporary listening. Τα σημερινά experiments δεν αναδρομολογούνται ως claims για ακροατές του 19ου αιώνα.

\begin{keybox}
\textbf{Το μοτίβο της καλής απάντησης:} αναγνώρισε το όριο, δείξε ποιο protocol το ελέγχει και περιόρισε το claim. Μην υπερασπίζεσαι μεγαλύτερο συμπέρασμα από αυτό που πραγματικά χρειάζεται η διατριβή.
\end{keybox}

\newpage
% PAGE 20
\section*{19. Καίρια ερωτήματα επιστημονικής ανατροφοδότησης — και τι να σημειώσεις}

\subsection*{Α. Scope και επιστημονική ταυτότητα}
\begin{enumerate}
\item «Αν έπρεπε να περιορίσουμε την πρόταση σε μία κύρια computational contribution και μία validation contribution, ποιες θα κρατούσατε;»
\item «Βλέπετε το PhD περισσότερο ως MIR representation/segmentation thesis ή ως computational musicology thesis με MIR εργαλεία;»
\item «Η M5 σας φαίνεται κατάλληλη ως τελική εφαρμογή ή θα προτιμούσατε να μείνει μικρότερη μέχρι να ωριμάσει η M1/M2;»
\end{enumerate}

\subsection*{Β. M1 και modelling}
\begin{enumerate}
\item «Θα αλλάζατε κάτι στο 3×4 \(U\times F\) πριν ξεκινήσει η υλοποίηση;»
\item «Θα επιμένατε σε semi-CRF από την αρχή ή θα βάζατε simpler segment model ως πρώτο implementation milestone;»
\item «Ποιο θα θεωρούσατε πειστικό primary criterion για representation selection;»
\end{enumerate}

\subsection*{Γ. M2/EISE}
\begin{enumerate}
\item «Σας πείθει ότι B2 και EISE πρέπει να διαφέρουν κυρίως στην provenance των relations;»
\item «Θα ζητούσατε stricter \(B2_{sim}\) από την αρχή;»
\item «Ποια diagnostic θα σας έπειθε ότι το identity projection δεν έχει καταρρεύσει και ότι τα local states κρατούν realization information;»
\end{enumerate}

\subsection*{Δ. Δεδομένα και επίβλεψη}
\begin{enumerate}
\item «Πόσο μεγάλο corpus θεωρείτε ρεαλιστικό για 4 χρόνια με σοβαρή επισημείωση;»
\item «Θα προτείνατε συν-επίβλεψη/στενή συνεργασία με βυζαντινομουσικολόγο για annotation και ιστορικές πηγές;»
\item «Ποια υποδομή/compute/repository discipline θα θέλατε να βλέπετε στους πρώτους 6 μήνες;»
\end{enumerate}

\subsection*{Ε. Δημοσιεύσεις και milestones}
\begin{enumerate}
\item «Ποιο πρώτο paper θα θεωρούσατε ισχυρότερο: symbolic representation benchmark ή EISE pilot;»
\item «Ποιο milestone θα θέλατε να έχει ολοκληρωθεί στο τέλος του πρώτου έτους;»
\end{enumerate}

\begin{warnbox}
\textbf{Σημείωσε όχι μόνο τι προτείνει, αλλά και τι θεωρεί επικίνδυνο.} Η πιο χρήσιμη πληροφορία μπορεί να είναι ποιο τμήμα θα κριθεί σκόπιμο να μικρύνει, ποιο θα θεωρήσει core και ποιο dataset θεωρεί αρκετό για δημοσίευση.
\end{warnbox}

\newpage
% PAGE 21
\section*{20. Τελική καρτέλα: 90 δευτερόλεπτα και 10 σύμβολα}
\subsection*{90-second pitch}
\begin{keybox}
«Η διατριβή ελέγχει αν η μουσική θέση μπορεί να αποτελέσει υπολογιστική μονάδα μεταβλητού μήκους, διατηρώντας τη μουσικολογική της σημασία. Στη M1 συγκρίνω 3 κλίμακες συμβολικής μονάδας και 4 σωρευτικά επίπεδα πληροφορίας και επιλέγω λειτουργική \(R^\star\) για την οριοθέτηση, υπό ελάχιστη απαίτηση οικογενειακής αναγνώρισης. Στη M2 συγκρίνω B2 και EISE: κοινή αρχιτεκτονική, αλλά διαφορετική προέλευση θετικών σχέσεων, απλή ενδοοικογενειακή έναντι ιστορικά τεκμηριωμένης εξήγησης. Με το ελάχιστο ανεξάρτητο εκτελεστικό κριτήριο της M4 ελέγχω αν η ιστορικά οργανωμένη παραλλαγή συμφωνεί με τη σημερινή εκτελεστική μεταβλητότητα. Η M3a επικυρώνει ανεξάρτητα το λεξιλόγιο, οι M3c/M3b εξετάζουν προσδοκία και κατηγοριοποίηση, και η M5 καταλήγει σε πολυδιάστατο χαρακτηρισμό πλήρων έργων του Νικολάου ως προς ρητό ιστορικό σώμα αναφοράς. Κάθε συμπέρασμα έχει καθορισμένα δεδομένα, ανεξάρτητη αξιολόγηση και σαφές όριο ισχυρισμού.»
\end{keybox}
\subsection*{Ο χάρτης 10 συμβόλων}
\begin{tabularx}{\linewidth}{>{\bfseries}p{2.3cm}X>{\bfseries}p{2.5cm}X}
\(R^\star\) & Επιλεγμένη αναπαράσταση & \(z_{id}\) & Οικογενειακή ταυτότητα\\
\(\mathcal R^\star\) & Αβεβαιότητα επιλογής & \(h_t\) & Τοπικές καταστάσεις\\
\(g_{min}\) & Ελάχιστη επάρκεια οικογενειών & \(\Delta_{hist}\) & EISE μείον B2\\
\(M(R)\) & Κριτήριο ορίων & \(I_F\) & Έκπληξη οικογένειας\\
\(G(R)\) & Οικογενειακό κριτήριο & \(D_{within}\) & Διαφορά πραγματώσεων\\
\end{tabularx}
\subsection*{Πριν από επιστημονική συζήτηση ή εξέταση}
Ξεχωρίζω κάθε φορά \textbf{(α)} τι έχει τεκμηριωθεί στις πηγές, \textbf{(β)} τι έχει οριστεί ως πρωτόκολλο, \textbf{(γ)} τι δεν έχει δοκιμαστεί ακόμη. Εξηγώ κάθε τεχνικό όρο ως «πρόβλημα \(\rightarrow\) μέθοδος \(\rightarrow\) είσοδος/έξοδος \(\rightarrow\) εξωτερικός έλεγχος». Δεν αποδίδω ιστορική καταγωγή σε απλή ομοιότητα και δεν ονομάζω την απόκλιση «λάθος». 
\begin{warnbox}
\textbf{Τρεις επιδιώξεις:} επιστημονική αξιολόγηση του πυρήνα· προτεραιότητες δεδομένων και πρώτου έτους· συμφωνία για τα απολύτως ελάχιστα παραδοτέα και τις υποστηριζόμενες επεκτάσεις.
\end{warnbox}
\newpage
\section*{21. Είσοδοι--έξοδοι που πρέπει να εξηγώ στον πίνακα}
\subsection*{Από παρτιτούρα σε εμφάνιση θέσης}
\begin{enumerate}
\item Μεταγραφή/εξαγωγή Νέας Μεθόδου σε σταθερή συμβολική ακολουθία με προέλευση ανά γεγονός.
\item Επιλογή \(R^\star\) από M1: \(u_1/u_2/u_3\) και \(N/T/M/L\).
\item Encoder: αριθμητική κατάσταση \(h_t\) ανά ατομικό γεγονός. Candidate spans: οποιαδήποτε επιλέξιμα τμήματα επιτρεπτού μήκους.
\item Segment pooling και \(z_{id}\): ταυτότητα οικογένειας· τα \(h_t\) διατηρούν πληροφορία τοπικής πραγμάτωσης.
\item Semi-Markov CRF: κοινή επιλογή ορίων και οικογενειών σε μεταβλητού μήκους τμήματα, με μερικές ετικέτες. Open-set: άγνωστη οικογένεια ή αποχή, όχι καταναγκαστική ένταξη.
\item Έξοδος μουσικολόγου: \((\mathrm{αρχή},\mathrm{τέλος},\mathrm{οικογένεια},\mathrm{αβεβαιότητα})\), μη καλυμμένες περιοχές και ανακτήσεις συγγενών πραγματώσεων με πηγές.
\end{enumerate}
\subsection*{Τι αλλάζει μεταξύ B2 και EISE;}
Κοινή αρχιτεκτονική, δεδομένα ανάπτυξης, οικογενειακές ετικέτες, budget και ισορροπημένες δυσκολίες. \textbf{Μόνο η πηγή εποπτείας της σχέσης} είναι η αποφασιστική μεταβολή: B2 ίδιο family χωρίς τεκμηριωμένη εξηγητική ακμή· EISE επαληθευμένα εξηγητικά ζεύγη. Ιστορική εξήγηση \emph{δεν} παράγεται συνθετικά.
\subsection*{Γιατί όχι απλή πρόβλεψη «νέας θέσης»;}
Άγνωστο οικογενειακό label σημαίνει ότι η εμφάνιση δεν μπορεί να αποδοθεί με ασφάλεια στο παρόν λεξιλόγιο, όχι ότι ανακαλύφθηκε αυτομάτως ιστορικά νέα θέση. Προστίθεται ανεξάρτητη μουσικολογική κρίση.
\subsection*{Προεκπαίδευση και συνθετικά δεδομένα}
Κοινή προεκπαίδευση των νευρωνικών συγκριτών σε συμβολικό ρεπερτόριο, ενδεχομένως με masked-event reconstruction· οι συνθετικές μεταβολές είναι \emph{ελεγχόμενα τεχνητά παραδείγματα}, όχι ιστορικοί μάρτυρες. Παράγωγα του ίδιου πρωτοτύπου μένουν στον ίδιο διαχωρισμό train/dev/test. Αριθμοί επιπέδων/κρυφών διαστάσεων είναι πιλοτικές επιλογές, όχι ήδη επιβεβαιωμένη βέλτιστη αρχιτεκτονική.
\say{«Είσοδος: ακολουθία συμβολικών γεγονότων. Έξοδος: μουσικολογικά ελέγξιμος χάρτης εμφανίσεων και οικογενειών με αβεβαιότητα και διατηρημένη πληροφορία πραγματώσεων.»}

\newpage
\section*{22. Έντεκα παγίδες από την προφορική προετοιμασία}
\begin{enumerate}
\item \textbf{«Ένα νεύμα = ένα γεγονός».} Όχι: \(u_1\) γραφικός χαρακτήρας, \(u_2\) ατομικό μελωδικό γεγονός.
\item \textbf{«Ο EISE βλέπει μόνο ακολουθίες γεγονότων».} Όχι: εξετάζει υποψήφια τμήματα, όρια και οικογένειες.
\item \textbf{«Ο EISE αναγνωρίζει ιστορική εξήγηση από τον ήχο».} Όχι: το ιστορικό τεκμήριο ανήκει στο μητρώο πηγών.
\item \textbf{«Άγνωστη οικογένεια = νέα θέση».} Όχι: μπορεί να είναι αβεβαιότητα, ανεπαρκές λεξιλόγιο ή λάθος.
\item \textbf{«Τυφλή M3a = οι ειδικοί δεν γνωρίζουν τίποτε».} Όχι: γνωρίζουν το αναγκαίο μουσικό συμφραζόμενο, όχι τις προβλέψεις/αμοιβαίες κρίσεις.
\item \textbf{«Η M3b επιλέγει απλώς τις πιο όμοιες μελωδίες».} Όχι: το ζητούμενο είναι η οικογενειακή σχέση \emph{πέρα από} την ομοιότητα.
\item \textbf{«Υψηλό \(I_F\) σημαίνει υψηλό \(I_{event}\)».} Όχι απαραίτητα: χρειάζονται αντίθετοι συνδυασμοί.
\item \textbf{«Το PPM αποδεικνύει την αξία του συμφραζομένου μόνο επειδή χρησιμοποιεί backoff».} Όχι: απαιτείται σύγκριση σε νέα έργα με \(P(F_t\mid M,L)\).
\item \textbf{«Η μεταβλητή \(R\) προσομοιώνει την πλήρη λατρεία».} Όχι: είναι μη μουσική ένδειξη συμφραζομένου.
\item \textbf{«Μια σύγχρονη εκτέλεση επαναφέρει συγκεκριμένη ιστορική εξήγηση».} Όχι χωρίς ανεξάρτητη πηγολογική τεκμηρίωση.
\item \textbf{«Η M5 μετρά την ποιότητα/κλασικότητα».} Όχι: μετρά πολυδιάστατες διαφοροποιήσεις ως προς ρητό σώμα αναφοράς.
\end{enumerate}
\subsection*{Ο κανόνας μίας προφορικής απάντησης}
\begin{keybox}
«Το μουσικολογικό πρόβλημα είναι [Χ]. Η υπολογιστική μέθοδος [Υ] το αντιμετωπίζει μέσω [είσοδοι--έξοδοι]. Η επιλογή ελέγχεται έναντι [baseline] στο [ανεξάρτητο κριτήριο]. Αν δεν επιβεβαιωθεί, ο ισχυρισμός περιορίζεται σε [ελάχιστο έγκυρο αποτέλεσμα].»
\end{keybox}
\subsection*{Οι τρεις πιο χρήσιμες ερωτήσεις που θα κάνω}
(1) Ποιο είναι το ρεαλιστικό ελάχιστο πειραματικό πεδίο του πρώτου έτους; (2) Ποιο είναι το πιο εύθραυστο σημείο της M2: προέλευση ιστορικών ζευγών ή ανεξάρτητο εκτελεστικό κριτήριο; (3) Να προηγηθεί απλούστερο segment model πριν από την τελική δοκιμή semi-Markov;
\end{document}
